\section{Discussion}\label{sec:discussion}

\subsection{Related work and what is new here}\label{sec:relatedwork}
The algebraic core of our loss is not new, and we do not present it as such. The
identity $1-r=\tfrac12\,\mathrm{MSE}(z_y,z_{\hat y})$ is the elementary fact that
the Pearson correlation equals the inner product of the $z$-scored vectors, so
minimizing the MSE between standardized variables is the same as maximizing $r$;
this has been used implicitly wherever a correlation or cosine objective is
optimized. Correlation-, concordance- and ranking-based training objectives also
already exist in both the statistical-genetics and machine-learning literatures:
Lin's concordance correlation \citep{lin1989} has long served as a
scale-and-offset-aware agreement criterion, \citet{blondel2015} reframed genomic
selection as a learning-to-rank problem with NDCG, and a recent convolutional
genomic-prediction architecture substitutes a family of $L_p$ correlation-style
losses for MSE \citep{pnngs2024} --- although, as we note in the introduction, it
mistakenly asserts that ``Pearson correlation cannot be set as the loss
function.'' Differentiable Spearman and other ranking surrogates are likewise
established in the broader ML literature. Our contribution is therefore
\emph{not} the loss or the identity. What we add is (i) a systematic,
calibration-aware and selection-metric-aware empirical characterization of
correlation-consistent training across \NUM{101} trait--datasets, in which
architecture --- not the loss --- emerges as the dominant moderator; and (ii) a
closed-form, rank-preserving post-hoc affine recalibration
(Proposition~\ref{prop:2}) framed explicitly as the remedy for the
scale-invariance that any correlation objective discards. The numerically stable
$\tfrac12\,\mathrm{MSE}(z_y,z_{\hat y})$ formulation is a usable training target
rather than a formal novelty.

\subsection{A Pearson loss is a useful default, not a universal winner}
Our central empirical message is deliberately unromantic. Training a neural
genomic predictor with the standardized-MSE Pearson loss instead of MSE does
move predictive ability in the right direction --- pooled across the
\NUM{251} loss--model cells the mean improvement is $\Delta r = +0.038$ (BCa 95\%
CI $[0.028,0.050]$, $p_{\mathrm{Holm}}=1.2\times10^{-8}$), with matching gains in
Spearman $\rho$ ($+0.037$, $[0.028,0.048]$) and in NDCG@10 ($+0.016$,
$[0.011,0.021]$) --- but it is not the best loss, and the gain is concentrated in
a subset of architectures. The concordance-correlation loss ($\Delta r = +0.044$,
$[0.034,0.056]$) and the hybrid loss ($+0.041$, $[0.031,0.052]$) both edge out the
pure Pearson loss on essentially every metric we examined; the three
correlation-aware objectives are statistically separable from MSE but only
marginally separable from one another (Figure~\ref{fig:loss}). The pooled
$n=\NUM{251}$ figure is, however, an \emph{unweighted} average over an
architecturally unbalanced set of cells, and this must be kept in view: the CNN
and MLP were each run on all \NUM{101} trait--datasets, but the Transformer was
run on only \NUM{49} trait--datasets spanning four species (lentil, maize, rice,
soybean), so $251=101+101+49$. Because the Transformer also carries the largest
per-architecture effect, the pooled mean is pulled upward by a smaller,
non-random subset of datasets, and per-architecture effects are not estimated on a
common dataset base. The stratified estimates make the heterogeneity explicit: the
gain is large for the Transformer ($\Delta r = +0.101$, $[0.071,0.131]$, on its
\NUM{49}-dataset subset) and the 1D-CNN ($+0.044$, $[0.026,0.066]$) yet vanishes
for the MLP ($+0.0027$, $[-0.0004,0.0055]$), whose confidence interval includes
zero and whose CCC effect is frankly null ($+0.0004$, $p_{\mathrm{Holm}}=0.17$).

The confirmatory mixed model is consistent with this picture but must be read for
what it actually estimates. It is additive in $\mathrm{loss}+\mathrm{model}+
\mathrm{scheme}$, with no $\mathrm{loss}\times\mathrm{model}$ interaction, so the
Pearson-loss fixed effect ($+0.039$, $[0.033,0.045]$) is an average loss effect
and the model terms are \emph{baseline} contrasts, not moderators. The
$\mathrm{model}_{\mathrm{MLP}}=+0.060$ coefficient means the MLP has the highest
intrinsic accuracy among the three networks (relative to the CNN reference), not
that architecture modulates the loss benefit; the genuine evidence for moderation
is the stratified per-model $\Delta$ above (Transformer $+0.101$ versus MLP
$+0.0027$), which is what we rely on when we say architecture dominates. The model
also reports very tight loss intervals ($t\approx12$--$15$); these should be read
with caution, because the random-effects structure is $(1\,|\,\mathrm{dataset})+
(1\,|\,\mathrm{dataset{:}trait})$ with no term for the ten within-cell folds
(\NUM{5} folds $\times$ \NUM{2} repeats), which reuse the same frozen architecture
and overlapping markers and are therefore not independent. The unstratified
per-cell BCa intervals above, which aggregate to one $\Delta$ per cell, are the
estimates we treat as primary.

The blunter reading is that none of the neural predictors, under any loss, was
competitive with the simple parametric baselines on these data. By mean rank
across all cells, ridge regression (rank $2.85$ on $r$, $4.01$ on NDCG@10) and
GBLUP (rank $3.45$, $4.32$) outrank \emph{every} neural-network/loss combination;
the best neural entry is the MLP with the hybrid loss (rank $5.28$ on $r$), and the
worst is the Transformer with MSE (rank $12.67$). These mean ranks compare
networks against baselines on \emph{absolute} performance, and the same HPO caveat
discussed below (the architecture is tuned on full-data validation splits) can
mildly inflate the neural numbers; that inflation works \emph{in favour} of the
networks and therefore only strengthens, rather than threatens, the conclusion
that the linear baselines win. The Pearson loss narrows the gap between a neural
network and a ridge model --- and on the deeper, more loss-sensitive
architectures it narrows it substantially --- but it does not close it. This is
consistent with the now-standard finding that neither deep learning nor GBLUP
systematically dominates in genomic prediction
\citep{montesinos2025dlgblup,montesinos2021power}: the loss function is one lever
among several, and on additive, marker-dense traits the linear baselines remain
hard to beat. The honest claim is therefore the conditional one: \emph{if} a
practitioner has already chosen a non-linear architecture --- particularly a
convolutional or attention-based one --- whose default MSE training is poorly
aligned with the correlation metric used for reporting and selection, then
switching to a correlation-consistent loss is a low-cost, reliably positive
change. It is not a reason to abandon GBLUP or ridge.

\subsection{Calibration is mandatory whenever phenotypic scale matters}
The flip side of Proposition~\ref{prop:1} is that a correlation loss is
\emph{blind} to scale and offset, and the raw outputs of a Pearson-trained
network confirm how badly this matters. On the Pearson-trained networks the raw
predictions carry a mean calibration slope of $4.25$ and a catastrophic mean
intercept of $-378.8$ (with enormous spread, intercept std $\approx5.1\times10^3$):
the predictions have the right \emph{shape} but live on the wrong axis entirely. A
single closed-form affine map (Proposition~\ref{prop:2}), fit on validation
predictions only, repairs this almost completely --- the mean bias collapses from
$-0.76$ to $+0.044$, the slope and intercept move toward their ideals, and RMSE
drops from $8.68$ to $7.07$ --- while leaving the Pearson correlation essentially
unchanged at the mean level ($0.531$ raw versus $0.533$ affine; the residual
difference is noise). The fitted affine slope $a^\star=\mathrm{cov}(y,p)/
\mathrm{var}(p)$ is estimated on validation data, so the rank-preservation
guarantee of Proposition~\ref{prop:2} holds only when $a^\star>0$. This is the
case for almost every cell (the mean affine slope is $+2.27$), but it is not
universal: in fewer than $1\%$ of folds ($90$ of $9{,}943$, $0.91\%$) a negative
validation covariance produces a decreasing map on test, flipping the sign of the
per-fold Pearson, with per-fold raw-versus-affine differences as large as $0.89$ in
the worst cases. We therefore qualify the theoretical statement: a strictly
increasing affine map ($a^\star>0$) preserves every correlation-type and
rank-based statistic exactly, and because $a^\star>0$ in over $99\%$ of folds the
correlation is preserved on average; the categorical phrasing that it
\emph{must} be preserved holds only conditionally on $a^\star>0$. Isotonic
calibration is marginally worse on $r$ ($0.516$). Isotonic regression is monotone
non-decreasing by construction (its mean fitted slope is $+0.76$), so the loss of
correlation is \emph{not} due to non-monotonicity; it arises because the isotonic
fit maps distinct inputs to identical plateau values (piecewise-constant ties),
which is the well-known nonlinearity-versus-rank trade-off of isotonic
regression. The practical consequence is unambiguous: a Pearson or CCC loss should
\emph{never} be deployed without recalibration when the downstream use requires
phenotypic values on their natural scale --- breeding values in trait units,
RMSE-based model comparison, or any decision that mixes predictions with measured
phenotypes (Figure~\ref{fig:calibration}). Affine calibration is the cheapest way
to recover the property the loss discards, and it costs two numbers.

That said, calibration is a corrective, not a source of accuracy, and our RMSE
results warn against expecting more from it than it can give. All RMSE deltas here
are computed on affine-calibrated predictions. Pooled, the Pearson loss does
\emph{not} reliably reduce calibrated RMSE relative to MSE: the point estimate is
$\Delta\mathrm{RMSE}=-0.010$ with a BCa CI that crosses zero ($[-0.055,0.079]$).
The paired Wilcoxon test is nominally significant ($p_{\mathrm{Holm}}=0.0049$), but
this is a sign/median statement, not a magnitude one --- the signed-rank test
detects a consistent direction in the per-cell differences while the mean effect
and its CI are indistinguishable from zero. We do not interpret this as an RMSE
improvement. The hybrid loss is similar (point estimate $+0.006$, CI
$[-0.041,0.097]$ spanning zero with a worse sign; $p_{\mathrm{Holm}}=1.5\times
10^{-4}$ again reflects only a weak directional regularity, not a meaningful change
in magnitude), and is in fact worse for the CNN ($+0.080$) while better only for
the Transformer ($-0.119$). Only the CCC loss --- which penalizes scale and offset
\emph{inside} the objective rather than after it --- delivers a consistent pooled
calibrated-RMSE reduction ($-0.057$, $[-0.133,-0.027]$, CI excluding zero), and
even that is driven by the CNN and Transformer. The confirmatory model agrees that
the Pearson-loss effect on calibrated RMSE is null ($-0.010$, $t=-0.17$, CI
spanning zero). In short, once predictions are correctly calibrated, the choice of
correlation loss buys ranking quality, not magnitude accuracy; if calibrated
magnitude is the explicit goal, the CCC loss is the better starting point.

\subsection{Selection quality requires top-$k$ and ranking metrics, and it is
noisier than $r$}
Because breeders act on the tails of the predicted distribution rather than on the
average fit, we evaluated selection directly. Here the correlation-consistent
losses again help on average --- relative efficiency at the top 10\% improves by
$\Delta = +0.038$ to $+0.040$ across the three losses, top-10\% overlap by $+0.014$
to $+0.018$, and NDCG@10 by $+0.016$ to $+0.017$ --- and the gains track the
predictive-ability gains, again concentrated in the CNN and Transformer and again
negligible for the MLP. For the MLP specifically the selection effects of the
Pearson loss are not distinguishable from zero ($\Delta$NDCG@10 $=+0.0026$,
$p_{\mathrm{Holm}}=0.28$; $\Delta$overlap@10 $=+0.0015$, $p_{\mathrm{Holm}}=1.0$;
both Holm-corrected, the convention we use throughout). Two cautions follow. First,
a better $r$ does not \emph{guarantee} better selection of the extremes: the
relationship between $\Delta r$ and $\Delta$NDCG@10 across cells is positive but
loose, exactly as \citet{blondel2015} and \citet{ornella2014} warned, which is
precisely why we report selection-aware metrics rather than inferring selection
quality from $r$ alone. Second, the most decision-relevant quantity is also the
noisiest: selection differential at $k=10\%$ has a positive pooled mean for the
Pearson loss ($+0.25$) but with an interval that, while excluding zero
($[0.019,0.612]$), is extremely wide because it depends on a handful of top
individuals per fold. NDCG@$k$ and relative efficiency are the better-behaved
summaries and are what we recommend reporting. The broader point reinforces
\citet{waldmann2019}: the metric used to choose and evaluate a model changes the
answer, and a genomic-selection study that reports only global $r$ is not
reporting the quantity it actually cares about.

\subsection{Robustness under difficult splits is, at best, suggestive}
We probed whether the benefit survives harder generalization regimes
(leave-family-out on SoyNAM, cross-environment transfer on CIMMYT wheat,
within-environment, and random $K$-fold) in a separate experiment. The honest
summary is that the gain weakens, and in places reverses, under hard splits, and
that the evidence base for these directional statements is thin and should be read
as suggestive rather than confirmatory. The leave-family-out, random and
within-environment contrasts each rest on only $n=4$ datasets with no $p$-values,
and the cross-environment contrast on $n=12$ \emph{ordered} environment pairs
derived from the same four environments on the same $599$ wheat lines, which are
not mutually independent; no multiple-testing correction is applied across the
roughly twenty split contrasts examined. With those caveats stated plainly: under
leave-family-out CV the Pearson-loss predictive gain is essentially absent
($\Delta r = -0.0030$, $[-0.0106,0.0010]$) and selection quality is, if anything,
slightly degraded ($\Delta$NDCG@10 $=-0.0042$, BCa CI $[-0.0089,-0.0005]$
excluding zero) --- but this CI is bootstrapped over four paired cells with no
$p$-value, so we read it as a suggestive signal of weakened benefit, not as
established significant degradation. Under cross-environment transfer the
predictive change is small and its inferential status is genuinely ambiguous: the
point estimate is $\Delta r=+0.0098$ with a BCa CI of $[0.0001,0.0204]$ that
narrowly excludes zero, yet the paired Wilcoxon test is non-significant
($p=0.151$). At $n=12$ the discrete signed-rank test and the bootstrap CI simply
disagree, and we do not resolve the disagreement in either direction; the
cross-environment selection metric shows no benefit ($\Delta$NDCG@10 $=-0.0012$,
$p=0.97$). The practically defensible reading is that the loss benefit is clearest
under random partitioning and is not demonstrated to carry over to family- or
environment-level extrapolation on the limited hard splits available here.

\subsection{Recommendations for breeders and practitioners}
We distil the above into concrete, defensible guidance.
\begin{enumerate}
\item \textbf{Match the loss to the architecture, not to fashion.} For linear
predictors (GBLUP, ridge), keep their native estimators --- they were the
strongest methods in this study and have no loss to tune. Reserve the
correlation-consistent loss for non-linear networks (CNN, Transformer), where it
gave the clearest and most reliable gains; for a plain MLP it is essentially
neutral and not worth the change.
\item \textbf{Always calibrate before reading predictions on their natural
scale.} If the deliverable is a ranking, raw correlation-trained outputs suffice;
if it is a breeding value, an RMSE, or anything compared against measured
phenotypes, apply the closed-form affine map fit on held-out data. It preserves
the ranking exactly whenever the fitted slope is positive (over $99\%$ of folds)
and is free.
\item \textbf{Report selection-aware metrics alongside $r$.} At minimum report
NDCG@$k$ and relative efficiency (or top-$k$ overlap) at the selection intensities
the program actually uses (here $k\in\{5,10,15,20\}\%$). Treat selection
differential as informative but high-variance.
\item \textbf{Prefer the CCC loss when calibrated scale and ranking both matter.}
Among the objectives tested, the CCC loss was the \emph{only} one to reduce pooled
calibrated RMSE (CI excluding zero) while matching the other correlation-aware
losses on ranking, giving a single training run that serves both magnitude and
selection use-cases. The hybrid loss did \emph{not} reliably improve pooled
calibrated RMSE (point estimate $+0.006$, CI crossing zero); it lowered RMSE only
for the Transformer and worsened it for the CNN, so its calibrated-RMSE benefit
should not be generalized.
\item \textbf{Do not expect a loss change to substitute for a better model.}
On additive, marker-dense traits the linear baselines remained ahead of all
neural variants; the loss is a refinement, not a replacement for sound model
choice.
\end{enumerate}

\subsection{Limitations and future work}
Several limitations temper these conclusions. First, and most importantly for
honest interpretation, our hyper-parameter search selected each network's
architecture on full-data validation splits under a neutral MSE objective, after
which the architecture was frozen and shared identically across all losses. Some
test-fold genotypes' phenotypes can therefore have informed the architecture
choice, so any \emph{absolute} accuracy we report --- raw and calibrated $r$
values, the calibration-table RMSE figures, the batch-size correlations, and the
neural rows of the mean-rank table --- may be mildly, and we stress symmetrically,
optimistically biased; none of these absolute numbers should be read as an
unbiased estimate of out-of-sample accuracy. Because every loss and every
calibration is compared \emph{within} the same frozen architecture, and our
inferential targets are always expressed as $\Delta$ versus MSE, this cannot bias
the relative loss/calibration comparisons that constitute our claims; GBLUP, having
no tuned hyper-parameters, is unaffected entirely, and the baseline-versus-network
ranking is if anything conservative because the bias favours the networks. A fully
nested HPO would remove even the cosmetic inflation of absolute numbers and is the
natural next step.

Second, the gains, while statistically supported on random splits, are modest in
absolute terms ($\Delta r\approx0.04$ pooled) and uneven across architectures and
species; they should not be oversold, and the pooled mean is computed over an
architecturally unbalanced set of cells (CNN and MLP on all \NUM{101}
trait--datasets, Transformer on only \NUM{49} across four species), with the
largest per-architecture effect attaching to the smaller Transformer subset.
Third, the most breeding-relevant metric, the selection differential, is estimated
with high variance from few top individuals per fold, and larger panels or pooled
multi-fold estimators would sharpen it. Fourth, our evidence on harder
generalization regimes is limited. We caution against over-reading the LMM:
its \emph{scheme} factor has only two levels, ``predefined'' (the ten EasyGeSe
species, their native five-fold CV partitions) and ``random'' ($K$-fold on SoyNAM
and CIMMYT wheat), so it is fully confounded with dataset source and the
$\beta_{\mathrm{random}}=+0.072$ ($[-0.151,0.296]$) coefficient contrasts those
two dataset groups, \emph{not} structured-versus-random splitting; moreover it is a
main effect on the absolute metric, not a loss$\times$scheme interaction, so a null
$\beta_{\mathrm{random}}$ says nothing about whether the loss benefit persists
across schemes. The leave-family-out and cross-environment contrasts are not levels
of this factor and are estimated only in the separate splits experiment, which ---
as discussed above --- suggests the benefit weakens or reverses under family- and
environment-level extrapolation rather than persisting. The number of genuinely
independent hard splits is small ($n=4$ families, $n=12$ non-independent
environment pairs), and establishing whether correlation-consistent training helps
under realistic structured prediction will require substantially more, and more
independent, hard-split data. Finally, our networks did not beat the linear
baselines, so the practical ceiling on the value of any loss here is bounded by the
predictors that benefit from it.

Future work follows directly. The group-aware variant of the loss
(Section~\ref{sec:groupaware}), which targets within-environment or within-family
correlation rather than the pooled value, is the most promising extension for
multi-environment breeding and was not exercised here. The mini-batch bias of the
correlation loss (Proposition~\ref{prop:3}) motivates principled large-batch or
statistic-accumulation training schemes. And a multi-trait formulation
$\mathcal{L}=\sum_t w_t(1-r_t)$ would let a single network serve correlated
breeding objectives jointly. In each case the same discipline applies: align the
training objective with the evaluation objective, calibrate when scale matters, and
judge selection by selection metrics.